\section*{Chapitre \ref{ch:g4:geometry} --- Droites et polygones}

\begin{solution}{exo:g4:geometry:1}
Constructions avec \cref{met:g4:geometry:draw}~; le petit carré
marque l'\omterm{def:g3:shapes:rightangle}{angle droit} entre $d$ et la \omterm{def:g4:geometry:def}{perpendiculaire}.
\end{solution}

\begin{solution}{exo:g4:geometry:2}
E~: les trois traits horizontaux sont \omterm{def:g4:geometry:def}{parallèles}~; chacun est
\omterm{def:g4:geometry:def}{perpendiculaire} au vertical. H~: les deux traits verticaux sont
\omterm{def:g4:geometry:def}{parallèles}, la barre transversale est \omterm{def:g4:geometry:def}{perpendiculaire} aux deux.
N~: les deux traits verticaux sont \omterm{def:g4:geometry:def}{parallèles}~; la diagonale
n'est ni \omterm{def:g4:geometry:def}{parallèle} ni \omterm{def:g4:geometry:def}{perpendiculaire}. T~: les deux traits sont
\omterm{def:g4:geometry:def}{perpendiculaires}. Z~: les traits du haut et du bas sont
\omterm{def:g4:geometry:def}{parallèles}~; la diagonale n'est ni l'un ni l'autre.
\end{solution}

\begin{solution}{exo:g4:geometry:3}
\omterm{def:g4:geometry:polygon}{Hexagone}~; \omterm{def:g4:geometry:polygon}{quadrilatère}~; triangle~; \omterm{def:g4:geometry:polygon}{octogone}~; \omterm{def:g4:geometry:polygon}{pentagone}. Un
\omterm{def:g4:geometry:polygon}{polygone} a autant de \omterm{def:g4:geometry:polygon}{sommets} que de côtés~: $6$, $4$, $3$, $8$,
$5$.
\end{solution}

\begin{solution}{exo:g4:geometry:4}
Un \omterm{def:g4:geometry:polygon}{pentagone} a $5$ côtés. Depuis $A$, les non-voisins sont $C$ et
$D$~: deux diagonales, $[AC]$ et $[AD]$.
\end{solution}

\begin{solution}{exo:g4:geometry:5}
$OA = 4$~cm (\omterm{def:g3:shapes:shapes}{rayon})~; $OM = 4$~cm (tout point du cercle)~;
$AB = 8$~cm (un \omterm{def:g4:geometry:circle}{diamètre} vaut deux fois le \omterm{def:g3:shapes:shapes}{rayon}).
\end{solution}

\begin{solution}{exo:g4:geometry:6}
Ouvrir le compas le long d'un côté, puis reporter cette ouverture
sur les deux autres côtés~: si l'écart pointe-crayon s'ajuste
exactement à un second côté, ces deux côtés sont égaux.
\end{solution}

\begin{solution}{exo:g4:geometry:7}
Le chemin se ferme en un \omterm{def:g4:geometry:polygon}{polygone} à $5$ côtés --- un \omterm{def:g4:geometry:polygon}{pentagone}
(pas régulier~!).
\end{solution}

\begin{solution}{exo:g4:geometry:8}
«~Les \omterm{def:g4:geometry:def}{droites perpendiculaires} se croisent toujours~»~: vrai ---
elles se croisent à l'\omterm{def:g3:shapes:rightangle}{angle droit} lui-même.

«~Les \omterm{def:g4:geometry:def}{droites parallèles} ne se croisent jamais~»~: vrai --- c'est
leur définition.

«~Une droite peut être \omterm{def:g4:geometry:def}{parallèle} à une droite et \omterm{def:g4:geometry:def}{perpendiculaire} à
une autre~»~: vrai --- sur papier quadrillé, une droite horizontale
est \omterm{def:g4:geometry:def}{parallèle} à toute droite horizontale et \omterm{def:g4:geometry:def}{perpendiculaire} à
toute droite verticale.
\end{solution}

\begin{solution}{exo:g4:geometry:9}
Les droites $d$ et $f$ semblent \omterm{def:g4:geometry:def}{parallèles} --- et elles le sont
toujours~: deux \omterm{def:g4:geometry:def}{droites perpendiculaires} à une même droite sont
\omterm{def:g4:geometry:def}{parallèles} (\cref{prop:g6:lines:facts} le rend officiel).
\end{solution}

\begin{solution}{exo:g4:geometry:10}
Un programme correct~: «~1. Tracer un carré $ABCD$ de côté $4$~cm.
2. Marquer le point $O$ où ses diagonales se croisent. 3. Tracer le
cercle de centre $O$ et de \omterm{def:g3:shapes:shapes}{rayon} $2$~cm.~» Le cercle touche chaque
côté en son milieu.
\end{solution}

\begin{solution}{exo:g4:geometry:11}
\omterm{def:g4:geometry:polygon}{Quadrilatère}~: $2$ diagonales. \omterm{def:g4:geometry:polygon}{Pentagone}~: $5$. \omterm{def:g4:geometry:polygon}{Hexagone}~: $9$.
Chaque nouveau \omterm{def:g4:geometry:polygon}{sommet} ajoute de plus en plus de diagonales --- les
comptes croissent de plus en plus vite ($+3$, puis $+4$, \dots).
\end{solution}
